\section{同步 vs 异步工具：AI 编程的三个时代}

\subsection{嘉宾讲座：Silas Alberti (Cognition/Devin)}

Cognition 研究负责人 Silas Alberti 分享了 AI 编程工具的演进路径和最佳使用策略。

\subsection{AI 编程工具的三个时代}

\begin{enumerate}
    \item \textbf{Code Completion}（$\sim$10\% 效率提升）：GitHub Copilot 加速编码
    \item \textbf{IDE Automation}（$\sim$20\% 效率提升）：AI IDE 实现单人任务完成
    \item \textbf{AI Software Engineer}（6--12$\times$ 效率提升）：云端 Agent 实现并行工作流
\end{enumerate}

\begin{importantbox}{从 10\% 到 6--12$\times$：量变到质变}
从 code completion 到 cloud agent，效率提升从百分比级跃升到倍数级。关键区别在于：cloud agent 实现了\textbf{并行}——一个开发者可以同时管理多个 Agent 执行不同任务。
\end{importantbox}

\subsection{同步 vs 异步工具}

\textbf{同步工具}（Sync）：
\begin{itemize}
    \item 单线程，人在环中（human-in-the-loop）
    \item 注意力聚焦于单一任务
    \item Agent 工作 20 秒到 1.5 分钟
    \item 代表：Cursor、Windsurf
\end{itemize}

\textbf{异步工具}（Async）：
\begin{itemize}
    \item 多线程，人类委派任务后切换注意力
    \item Agent 工作 10 分钟到数小时
    \item 代表：Devin
\end{itemize}

\begin{warningbox}{``Semi-Async'' 陷阱}
在 30 秒到 3 分钟的``半异步''区间要特别小心——它太慢以至于无法保持心流（flow），又太短以至于不值得切换到其他任务。正确策略是：要么让工具更快以保持心流，要么投入更多时间换取更高的智能水平。
\end{warningbox}

\subsection{2025年的最佳编程工作流}

Silas 提出的最佳实践将三个阶段分配给不同的工具：

\begin{itemize}
    \item \textbf{Planning}（规划）：同步工具——使用 DeepWiki、Codemaps 理解代码库
    \item \textbf{Coding}（编码）：异步工具——委派给 Devin 执行
    \item \textbf{Testing}（测试）：同步工具——在 Windsurf 中测试和微调 Agent 的结果
\end{itemize}

\subsection{未来展望}

\begin{enumerate}
    \item 人类工程师成为 Agent 管理者：同步工具解决最困难的问题，异步工具获取 10$\times$ 杠杆
    \item 未来的核心技能：\textbf{委派与多线程}、\textbf{代码阅读}、\textbf{规划、范围界定和架构设计}
    \item 长期趋势：规划、编码、测试三个阶段都将逐步实现异步化和自主化
\end{enumerate}

\subsection{团队采用 AI IDE 的评估框架}

如果把 Week 3 的内容转成组织层面的判断标准，可以用下面三个问题快速评估一套 AI IDE 是否值得推广：

\begin{itemize}
    \item \textbf{它是否真正减少上下文切换？} 如果开发者仍要在大量外部文档和 prompt 模板间来回跳转，工具并未带来本质改进。
    \item \textbf{它是否让代码库更透明？} 好的 AI IDE 会放大 repo 中已经存在的结构化信息，而不是掩盖混乱。
    \item \textbf{它是否有稳定的回退机制？} 当 Agent 给出错误修改时，团队能否依赖测试、review 和清晰的 diff 快速恢复？
\end{itemize}

\begin{knowledgebox}{真正可落地的 AI IDE 必须同时优化人和仓库}
AI IDE 从来不是单纯的模型问题。它同时要求更好的代码组织、更稳定的评测体系和更清晰的项目说明。也正因此，context engineering 最终会倒逼团队重构自己的知识管理方式。
\end{knowledgebox}

\subsection{本章小结}

AI 编程工具从 code completion 演进到 cloud agent，带来了从 10\% 到 6--12$\times$ 的效率飞跃。关键是掌握同步与异步工具的正确使用场景，避免``半异步''陷阱，将规划-编码-测试分配给最合适的工具类型。

